\subsubsection{Conclusiones e implicaciones}

El análisis realizado permite extraer una serie de conclusiones claras sobre el comportamiento temporal del sistema.

En primer lugar, el pipeline de control (decisión, arbitraje y ejecución) presenta un comportamiento estable y acotado, tanto en términos de magnitud como de variabilidad. La latencia física se mantiene en un rango estrecho, con baja dispersión y sin presencia de eventos extremos, lo que indica un funcionamiento determinista y predecible de estos subsistemas.

En contraste, el módulo de percepción introduce un cambio cualitativo en el comportamiento del sistema. No solo incrementa significativamente la latencia end-to-end, sino que además actúa como principal fuente de variabilidad. La similitud entre la dispersión de la latencia de percepción y la latencia end-to-end percibida, tanto en desviación estándar como en percentiles altos, evidencia que este subsistema domina el comportamiento global del sistema.

Adicionalmente, la presencia de colas largas y eventos de alta latencia en la percepción indica la existencia de episodios de degradación temporal que no pueden explicarse como ruido, sino como un comportamiento estructural asociado a la naturaleza asíncrona y dependiente de los datos del proceso de localización.

Desde una perspectiva de sistema, estos resultados implican que el rendimiento no está limitado por la capacidad de decisión o ejecución, sino por el cierre del bucle de percepción. En consecuencia, el sistema opera en un régimen \textit{perception-bound}, donde las mejoras en el pipeline de control tienen un impacto limitado sobre la latencia total.

Por tanto, cualquier estrategia orientada a reducir la latencia o mejorar la previsibilidad del sistema debe centrarse prioritariamente en el módulo de percepción. Esto incluye la optimización de la frecuencia de actualización o la reducción del coste computacional de los algoritmos de localización.

Finalmente, los resultados ponen de manifiesto que la latencia percibida no solo es mayor que la latencia física, sino que presenta una naturaleza estadística diferente, caracterizada por mayor dispersión, asimetría y presencia de valores extremos. Esta diferencia debe ser tenida en cuenta en el diseño y evaluación de sistemas robóticos, ya que afecta directamente a la capacidad de respuesta y a la estabilidad temporal del sistema en condiciones reales de operación.